\chapter{Estimating the share from repeated measurements and counts}
\label{estimation:chapter}

The theoretical share requires a breeding-value distribution and a conditional
fitness relationship. Neither is directly observed in the worked example.
Instead, each genotype has repeated trait measurements and one offspring
count. This chapter derives the estimator from that observation model. It
then derives predictions on the observed-data scale, the specified shrinkage
comparison, and the quartic calculation that separates a local approximation
from variation in a complete curve.

The probability calculations are exact under the stated model. Parameter
estimates are obtained by numerical optimisation, and bootstrap intervals
describe a fitted sampling model. Deriving a likelihood does not prove that its
assumptions hold in a biological dataset, that a numerical optimiser finds a
global optimum, or that a particular interval has its nominal coverage.

\section{Specify the observations and the latent quantity}
\label{estimation:model}

Let \(A_i\) be the breeding value of genotype \(i\), let \(\mu\) be the
population mean, and write \(X_i=A_i-\mu\). There are \(n\) unrelated genotypes
and \(m\) independent trait measurements per genotype. The teaching model is
\begin{align}
 X_i&\sim N(0,G),\nonumber\\
 Y_{ij}&=\mu+X_i+e_{ij},
 & e_{ij}&\sim N(0,V_E),\qquad j=1,\ldots,m.
 \label{estimation:measurement_model}
\end{align}
The \(X_i\) are independent across genotypes, and all measurement errors are
independent of them and of one another. Here \(V_E\) is measurement-environment
variance. In this example it is not an additional environmental variance to be
integrated into an underlying phenotype-fitness function.

Let \(R_i\) denote the offspring count, with observed value \(r_i\). It is the
instance of reproductive output \(W\) in this observation model. A separate
assay gives
\begin{align}
 R_i\mid X_i=x&\sim\operatorname{Poisson}\{\lambda(x)\},\nonumber\\
 \lambda(x)&=W_A(\mu+x)
 =\exp\{\alpha+bx+\tfrac12Hx^2\}.
 \label{estimation:count_model}
\end{align}
Conditional on \(X_i\), reproduction is independent of the environments in the
trait measurements. The parameter \(\alpha=\ell_A(\mu)\) is log expected output
at the population mean. Thus \(e^\alpha\) is not generally the population mean
output \(\E[\lambda(X)]\).

The generating example has \(n=800\), \(m=3\), \(G=V_E=1\),
\(\mu=0\), \(\alpha=\log4\), \(b=0.7\), and \(H=-0.7\). The trait has centimetre
units, \(G,V_E\) have \({\rm cm}^2\) units, \(b\) has \({\rm cm}^{-1}\) units,
and \(H\) has \({\rm cm}^{-2}\) units. The true breeding values are saved
separately and are not supplied to the estimator.

\paragraph{Biological interpretation.}
Repeated measurements distinguish variation within genotypes from variation
between them. Calling the latter additive genetic variance requires the
additional assumption that the genetic contribution to this trait is wholly
additive. Permanent environmental differences, dominance, epistasis or a
shared environment with the reproductive assay would change the model.
Clonal replication alone does not eliminate those possibilities. The example
isolates the estimation problem under explicit assumptions; it is not a
general identification argument for clonal data.

\section{Calibrate the measurement model}
\label{estimation:calibration}

Define the genotype mean and the mean across genotypes by
\[
 \bar Y_i=\frac1m\sum_{j=1}^mY_{ij},
 \qquad
 \bar Y=\frac1n\sum_{i=1}^n\bar Y_i.
\]
Because \(\bar Y_i=\mu+X_i+\bar e_i\), independence gives
\begin{equation}
 \E[\bar Y_i]=\mu,\qquad
 \Var(\bar Y_i)=G+\frac{V_E}{m}.
 \label{estimation:mean_measurement_variance}
\end{equation}
The factor \(1/m\) follows from
\(\Var(m^{-1}\sum_j e_{ij})=m^{-2}\sum_jV_E=V_E/m\).
It is not obtained if the repeated environments share an unmodelled effect.

\begin{proposition}[Balanced variance calibration]
\label{estimation:anova_calibration}
Under Equation~\eqref{estimation:measurement_model}, with \(n\geq2\) and
\(m\geq2\), the estimators
\begin{align}
 \widehat V_E
 &=\frac{\sum_{i=1}^n\sum_{j=1}^m(Y_{ij}-\bar Y_i)^2}{n(m-1)},\nonumber\\
 \widehat G
 &=\frac{\sum_{i=1}^n(\bar Y_i-\bar Y)^2}{n-1}
   -\frac{\widehat V_E}{m},
 \qquad
 \widehat\mu=\bar Y
 \label{estimation:calibration_estimators}
\end{align}
are unbiased for \(V_E,G,\mu\), respectively.
\end{proposition}

\begin{proof}
Within a genotype, subtracting its mean removes both \(\mu\) and \(X_i\):
\[
 Y_{ij}-\bar Y_i=e_{ij}-\bar e_i.
\]
Expanding the residual sum of squares gives
\begin{align*}
 \sum_{j=1}^m(e_{ij}-\bar e_i)^2
 &=\sum_j e_{ij}^2-2\bar e_i\sum_j e_{ij}+m\bar e_i^2\\
 &=\sum_j e_{ij}^2-m\bar e_i^2.
\end{align*}
Its expectation is \(mV_E-m(V_E/m)=(m-1)V_E\). Summing over \(n\)
genotypes and dividing by \(n(m-1)\) proves the first claim.

For the between-genotype term, put \(T=G+V_E/m\). The identity
\[
 \sum_i(\bar Y_i-\bar Y)^2
 =\sum_i(\bar Y_i-\mu)^2-n(\bar Y-\mu)^2
\]
has expectation \(nT-n(T/n)=(n-1)T\), because the genotype means are
independent and \(\Var(\bar Y)=T/n\). Consequently their sample variance
estimates \(G+V_E/m\) without bias. Subtracting \(\widehat V_E/m\) gives
\[
 \E[\widehat G]
 =G+\frac{V_E}{m}-\frac{\E[\widehat V_E]}{m}=G.
\]
This expectation argument does not require the two estimated terms to be
independent. Finally, averaging the genotype means gives
\(\E[\widehat\mu]=\mu\).
\end{proof}

Unbiasedness does not ensure that every realised estimate of \(G\) is positive.
When within-genotype noise is large relative to the between-genotype variation,
the subtraction can be zero or negative. The supplied implementation rejects
calibrations with \(\widehat G\leq10^{-8}\); it does not replace them with a
small positive variance. This is a stated numerical admissibility rule,
not a proof that smaller true variances cannot occur.

\section{Derive the distribution of an uncertain breeding value}
\label{estimation:conditioning}

For the moment, regard \(\mu,G,V_E\) as known, with \(G,V_E>0\). Write
\(d_i=\bar y_i-\mu\) for the centred mean of the observed measurements.
The normal prior density for \(X_i=x\) is proportional to
\(\exp(-x^2/(2G))\). The conditional measurement density is proportional to
\[
 \exp\!\left\{-\frac{1}{2V_E}
 \sum_j(y_{ij}-\mu-x)^2\right\}.
\]
Separate the within-genotype variation from its mean:
\begin{align}
 \sum_j(y_{ij}-\mu-x)^2
 &=\sum_j\{(y_{ij}-\bar y_i)+(d_i-x)\}^2\nonumber\\
 &=\sum_j(y_{ij}-\bar y_i)^2+m(d_i-x)^2.
 \label{estimation:measurement_factorisation}
\end{align}
The cross term vanishes because
\(\sum_j(y_{ij}-\bar y_i)=0\). The first term does not depend on \(x\).
Thus the genotype mean contains all the information about \(X_i\) supplied
by these balanced measurements when the variance parameters are fixed.

\begin{proposition}[Normal posterior and prediction reliability]
\label{estimation:posterior}
Define
\[
 R_{\rm pred}=\frac{G}{G+V_E/m},\qquad
 M_i=R_{\rm pred}(\bar y_i-\mu),\qquad
 v=G(1-R_{\rm pred}).
\]
Then
\begin{equation}
 X_i\mid Y_{i1}=y_{i1},\ldots,Y_{im}=y_{im}
 \sim N(M_i,v).
 \label{estimation:posterior_formula}
\end{equation}
Across populations generated by the model,
\begin{equation}
 \Var(M)=R_{\rm pred}G,\qquad
 \Cov(X,M)=R_{\rm pred}G,\qquad
 \Corr(X,M)^2=R_{\rm pred}.
 \label{estimation:reliability_identities}
\end{equation}
\end{proposition}

\begin{proof}
Multiplying prior and measurement densities gives a kernel with exponent
\[
 -\frac12\left\{\frac{x^2}{G}
      +\frac{(x-d_i)^2}{V_E/m}\right\}.
\]
Write \(T_E=V_E/m\). Expanding and completing the square gives
\begin{align*}
 \frac{x^2}{G}+\frac{(x-d_i)^2}{T_E}
 &=\left(\frac1G+\frac1{T_E}\right)x^2
   -\frac{2d_i}{T_E}x+\frac{d_i^2}{T_E}\\
 &=\frac{(x-M_i)^2}{v}
   +\frac{d_i^2}{T_E}-\frac{M_i^2}{v},
\end{align*}
where
\[
 v=\left(\frac1G+\frac1{T_E}\right)^{-1}
   =\frac{GT_E}{G+T_E},
 \qquad
 M_i=v\frac{d_i}{T_E}
   =\frac{G}{G+T_E}d_i.
\]
The terms not involving \(x\) are absorbed into the normalising constant.
The remaining density is normal with the stated mean and variance.

For the prediction identities, \(M=R_{\rm pred}(X+\bar e)\), with
\(\Var(X+\bar e)=G+T_E\). Therefore
\[
 \Var(M)=R_{\rm pred}^2(G+T_E)=R_{\rm pred}G.
\]
Independence of \(X\) and \(\bar e\) gives
\[
 \Cov(X,M)=R_{\rm pred}\Cov(X,X+\bar e)=R_{\rm pred}G.
\]
Squaring the correlation and substituting both variances gives
\[
 \Corr(X,M)^2
 =\frac{(R_{\rm pred}G)^2}{G(R_{\rm pred}G)}
 =R_{\rm pred}.
\]
These identities require positive variance for the correlation to be defined.
\end{proof}

The law of total variance makes the meaning of shrinkage explicit:
\begin{align}
 G=\Var(X)
 &=\Var\{\E[X\mid Y_1,\ldots,Y_m]\}
   +\E\{\Var(X\mid Y_1,\ldots,Y_m)\}\nonumber\\
 &=R_{\rm pred}G+v.
 \label{estimation:total_variance}
\end{align}
The variance among predictions is smaller than the biological variance because
the predictions omit unresolved variation. It is not evidence that the
population itself has contracted. As \(V_E/m\) tends to zero,
\(R_{\rm pred}\) tends to one and \(v\) tends to zero. As measurement noise
increases without bound, the predictions tend to zero while uncertainty tends
to the prior variance \(G\).

The implementation substitutes \(\widehat\mu,\widehat G,\widehat V_E\) into
these expressions. The resulting normal distributions are conditional on
calibrated parameter values. They are not full Bayesian posteriors that also
integrate over uncertainty in those values. Repeating calibration in the
bootstrap addresses their sampling uncertainty under the specified model.
Using predictions as error-free covariates discards the unresolved variance,
a concern discussed by \citet{HadfieldEtAl2010}.

\section{Construct the likelihood for offspring counts}
\label{estimation:likelihood}

Write \(\theta=(\alpha,b,H)^{\mathsf T}\) and
\[
 \eta_\theta(x)=\alpha+bx+\tfrac12Hx^2,\qquad
 \lambda_\theta(x)=e^{\eta_\theta(x)}.
\]
Given the calibrated measurements, the count probability for genotype \(i\)
is obtained by averaging over its uncertain breeding value:
\begin{align}
 L_i(\theta)
 &=P(R_i=r_i\mid y_{i1},\ldots,y_{im};\theta)\nonumber\\
 &=\int_{-\infty}^{\infty}
   \frac{\lambda_\theta(x)^{r_i}e^{-\lambda_\theta(x)}}{r_i!}
   \phi(x;M_i,v)\,dx\nonumber\\
 &=\int
   \exp\{r_i\eta_\theta(x)-e^{\eta_\theta(x)}
                  -\log(r_i!)\}\phi(x;M_i,v)\,dx.
 \label{estimation:mixture_likelihood}
\end{align}
Here \(\phi(x;M_i,v)\) is a normal density with variance \(v\). The log likelihood
is \(\ell(\theta)=\sum_i\log L_i(\theta)\). The logarithms are of
probabilities; no logarithm of \(r_i\) is taken.

For a zero count the Poisson factor is \(e^{-\lambda_\theta(x)}\), since
\(0!=1\). Zero counts therefore contribute ordinary, finite probabilities.
For \(v=0\), the conditional distribution of \(X_i\) is concentrated at
\(M_i\), so
\[
 L_i(\theta)
 =\frac{\lambda_\theta(M_i)^{r_i}
           e^{-\lambda_\theta(M_i)}}{r_i!}.
\]
That is the known-covariate limit. It is not the general case when \(v>0\).
In particular, averaging a Poisson probability is not the same operation as
evaluating it at the posterior mean or at the average rate.

This is a two-stage estimator. The measurement likelihood supplies the
calibration; the count likelihood then estimates the fitness coefficients.
Conditional on fixed calibration, genotypes contribute independent factors
\(L_i\). A joint analysis would also estimate the measurement parameters using
all components of a joint likelihood. The example does not claim equivalent
efficiency to such an analysis.

\section{From a normal integral to an evaluated likelihood}
\label{estimation:quadrature}

\subsection{The Hermite convention and normalisation}

Substitute \(x=M_i+\sqrt v\,t\) into
Equation~\eqref{estimation:mixture_likelihood}. Since
\[
 \phi(x;M_i,v)\,dx
 =\frac{1}{\sqrt{2\pi}}e^{-t^2/2}\,dt,
\]
the integral is a standard normal expectation. Let \(t_j\) and \(u_j\),
\(j=1,\ldots,q\), be Gaussian quadrature nodes and weights for the weight
function \(e^{-t^2/2}\). Define \(w_j=u_j/\sqrt{2\pi}\). Then
\begin{equation}
 L_i(\theta)\ \simeq\
 \sum_{j=1}^q w_j
 \exp\{r_i\eta_\theta(x_{ij})-e^{\eta_\theta(x_{ij})}
                     -\log(r_i!)\},
 \quad x_{ij}=M_i+\sqrt v\,t_j.
 \label{estimation:quadrature_likelihood}
\end{equation}
The weights \(w_j\) are positive and sum to one. This is the convention used
by \codepath{roots_hermitenorm}. It differs from the convention with weight
\(e^{-t^2}\), for which both the scaling of the nodes and the normalising
constant differ.

\begin{proposition}[Why the normal quadrature rule integrates polynomials]
\label{estimation:hermite_rule}
The \(q\)-node rule above integrates every polynomial of degree at most
\(2q-1\) exactly under the standard normal distribution.
\end{proposition}

\begin{proof}
Let \({\rm He}_q\) be the degree-\(q\) Hermite polynomial associated with the
standard normal weight. Its defining derivative formula is
\[
 {\rm He}_q(t)=(-1)^q e^{t^2/2}
             \frac{d^q}{dt^q}e^{-t^2/2}.
\]
If \(a\) is a polynomial of degree below \(q\), integration by parts \(q\)
times gives \(\int a(t){\rm He}_q(t)\phi(t;0,1)\,dt=0\). The boundary
terms vanish because a polynomial times a Gaussian density and its
derivatives tends to zero at both endpoints; the final derivative of \(a\)
is zero.

An orthogonal degree-\(q\) polynomial has \(q\) distinct real roots. To see
why, multiply its factors corresponding to real sign changes. If there were
fewer than \(q\), this multiplier would have degree below \(q\), while its
product with \({\rm He}_q\) would have one nonzero sign on the real line
except at isolated roots. Its Gaussian integral could not be zero, in
contradiction to orthogonality. Thus the \(q\) nodes can be chosen as those
distinct roots.

For a polynomial \(p\) of degree at most \(2q-1\), polynomial division gives
\(p=a\,{\rm He}_q+r\), where both \(a\) and \(r\) have degree at most
\(q-1\). Orthogonality removes the first term from the integral. At each
node, \({\rm He}_q(t_j)=0\), so \(p(t_j)=r(t_j)\). Interpolate \(r\) using
the Lagrange polynomials \(l_j\) through the nodes:
\[
 r(t)=\sum_jr(t_j)l_j(t),\qquad
 w_j=\int l_j(t)\phi(t;0,1)\,dt.
\]
Integrating gives \(\int p\phi=\sum_jw_jp(t_j)\), proving exactness.
The constant polynomial gives \(\sum_jw_j=1\). Applying exactness to
\(l_j^2\), of degree \(2q-2\), also gives
\[
 w_j=\sum_kw_k l_j(t_k)^2
     =\int l_j(t)^2\phi(t;0,1)\,dt>0.
\]
\end{proof}

The general Gaussian quadrature construction is described in
\citet{NISTDLMF}. Exactness for polynomials explains checks such as
\(\sum_jw_jt_j^2=1\) and \(\sum_jw_jt_j^4=3\). The Poisson-mixture integrand
is not a polynomial, so this theorem does not make a finite-order likelihood
exact. Its numerical accuracy must be assessed separately.

\subsection{Evaluate sums without losing numerical range}

For a particular genotype, set
\[
 a_{ij}=\log w_j+r_i\eta_\theta(x_{ij})
          -e^{\eta_\theta(x_{ij})}-\log(r_i!).
\]
Direct exponentiation can underflow when every probability is small. If
\(c_i=\max_j a_{ij}\), the identity
\begin{equation}
 \log\sum_j e^{a_{ij}}
 =c_i+\log\sum_j e^{a_{ij}-c_i}
 \label{estimation:logsumexp}
\end{equation}
evaluates the same expression with all exponentials at most one. The
implementation uses \codepath{logsumexp}, and evaluates
\(\log(r_i!)\) as \(\log\Gamma(r_i+1)\).

\subsection{Derive the score used by the optimiser}
\label{estimation:score}

Let
\[
 \vect d(x)=(1,x,x^2/2)^{\mathsf T},
 \qquad
 \eta_\theta(x)=\vect d(x)^{\mathsf T}\theta.
\]
The derivative of one log Poisson probability is
\begin{align*}
 \nabla_\theta\{r_i\eta_\theta(x)-e^{\eta_\theta(x)}
                                -\log(r_i!)\}
 &=r_i\vect d(x)-e^{\eta_\theta(x)}\vect d(x)\\
 &=(r_i-\lambda_\theta(x))\vect d(x).
\end{align*}
For the finite quadrature sum, differentiating its logarithm therefore gives
\begin{equation}
 \nabla_\theta\log L_i
 \simeq\sum_j\pi_{ij}
        (r_i-\lambda_\theta(x_{ij}))\vect d(x_{ij}),
 \qquad
 \pi_{ij}=\frac{e^{a_{ij}}}{\sum_k e^{a_{ik}}}.
 \label{estimation:mixture_score}
\end{equation}
The same identity holds for the exact integral with a conditional expectation
in place of the sum when differentiation under the integral is justified.
For the finite implemented sum the derivative is an ordinary derivative of a
finite positive expression.

There are two different conditional distributions in this calculation.
The normal \(N(M_i,v)\) describes breeding value after observing the heights.
The weights \(\pi_{ij}\) additionally use the count likelihood, and approximate
the breeding-value distribution after observing both heights and reproduction.
They sum to one, but are generally not the original normal quadrature weights.
The implementation sums Equation~\eqref{estimation:mixture_score} over
genotypes and minimises the negative mean log likelihood; dividing by \(n\)
does not change its minimiser.

\subsection{Numerical checks are part of the estimator}

The primary model uses 41 quadrature nodes initially and constrains
\(H\leq0\). Its numerical coefficient bounds are
\[
 -7\leq\alpha\leq7,\qquad -5\leq b\leq5,\qquad -6\leq H\leq0.
\]
These are bounds in the example's stated units. They are not universal
biological limits. Multiple starting values reduce dependence on a particular
initial guess, and a fit near a bound is flagged. Neither step proves global
optimality of a marginal likelihood.

The integration check refits at increasing orders, including 81, 161 and
321. For consecutive orders, it requires both
\begin{align}
 \max_k\frac{|\theta_{{\rm new},k}-\theta_{{\rm old},k}|}
                   {1+|\theta_{{\rm new},k}|}&<10^{-4},\nonumber\\
 \frac{|\ell_{\rm new}-\ell_{\rm old}|}{n}&<10^{-6}.
 \label{estimation:integration_check}
\end{align}
When the comparison passes, the lower-order fit is returned with the
higher check order recorded. An unresolved comparison raises an error.
For the recorded default data, comparing 41 and 81 nodes changes \(b,H,\Ag\)
by less than \(4\times10^{-13}\). That is a reproducible numerical result for
this dataset, not an analytic error bound for every possible observation.
An independent adaptive integral and a closed mean formula provide additional
checks in the supplied tests.

\section{From fitted coefficients to the variance share}
\label{estimation:components}

The fitted log conditional mean is
\[
 \widehat\ell_A(\widehat\mu+x)
 =\widehat\alpha+\widehat b x+\tfrac12\widehat Hx^2.
\]
Its derivatives are
\[
 \frac{d\widehat\ell_A}{dx}=\widehat b+\widehat Hx,\qquad
 \frac{d^2\widehat\ell_A}{dx^2}=\widehat H.
\]
Thus the fitted coefficients give slope and curvature at the estimated
population mean, \(x=0\). If software instead fits \(cx^2\), the second
derivative is \(2c\), so \(H=2c\). The design in the supplied code uses
\(x^2/2\) and requires no subsequent doubling.

Under the Gaussian quadratic model, the variance partition derived in
Proposition~\ref{gaussian:variance_partition} gives the plug-in quantities
\begin{equation}
 \widehat\Vlin=\widehat b^{\,2}\widehat G,\qquad
 \widehat\Vquad=\tfrac12\widehat H^{\,2}\widehat G^{\,2},\qquad
 \widehat\Ag=\frac{\widehat\Vlin}
                  {\widehat\Vlin+\widehat\Vquad}.
 \label{estimation:component_estimates}
\end{equation}
They describe variation in log expected output across breeding values.
Poisson variation in individual counts is not an extra term in their
denominator. When both components vanish, the share is undefined.

For the generating parameters, the calculation is exact:
\[
 \Vlin=0.7^2(1)=0.49,\qquad
 \Vquad=\tfrac12(-0.7)^2(1)^2=0.245,\qquad
 \Ag=\frac{0.49}{0.735}=\frac23.
\]
The fitted default value, approximately \(0.610744\), is a result of maximising
the likelihood for the recorded simulated observations. There is no
closed-form algebraic argument that a finite dataset must return that value
or the generating value. Reporting both components alongside the ratio also
shows whether a similar share is supported by a similar amount of variation.

\section{Predict quantities on the observed-data scale}
\label{estimation:predictions}

A graph of counts against observed heights asks for the count distribution
given those measurements. It does not ask for the count distribution at an
exact breeding value. Fix a genotype's measurement information and write
\(X\mid Y\sim N(M,v)\). In this section parameters are treated as fixed.

\begin{proposition}[Predictive moments and zero counts]
\label{estimation:predictive_moments}
For a Poisson mixture with conditional rate \(\lambda(X)\),
\begin{align}
 \mu_R(Y)&=\E[R\mid Y]=\E[\lambda(X)\mid Y],\nonumber\\
 s_R^2(Y)&=\Var(R\mid Y)
   =\E[\lambda(X)\mid Y]+\Var(\lambda(X)\mid Y),\nonumber\\
 p_0(Y)&=P(R=0\mid Y)=\E[e^{-\lambda(X)}\mid Y].
 \label{estimation:predictive_formula}
\end{align}
The first two identities assume the indicated moments are finite.
\end{proposition}

\begin{proof}
Conditional on \(X\), a Poisson variable has mean and variance \(\lambda(X)\).
Conditional independence of reproduction and the measurement environments
therefore gives
\[
 \E[R\mid Y]=\E\{\E[R\mid X,Y]\mid Y\}
           =\E[\lambda(X)\mid Y].
\]
The law of total variance, conditional on \(Y\), gives
\begin{align*}
 \Var(R\mid Y)
 &=\E\{\Var(R\mid X,Y)\mid Y\}
   +\Var\{\E[R\mid X,Y]\mid Y\}\\
 &=\E[\lambda(X)\mid Y]+\Var(\lambda(X)\mid Y).
\end{align*}
Finally \(P(R=0\mid X,Y)=e^{-\lambda(X)}\); averaging that probability over
\(X\mid Y\) proves the third identity.
\end{proof}

The supplied routine evaluates these as weighted normal expectations. If
\(\lambda_j=\lambda(M+\sqrt v\,t_j)\), it returns
\[
 \mu_R\simeq\sum_jw_j\lambda_j,\qquad
 s_R^2\simeq\mu_R+\sum_jw_j\lambda_j^2-\mu_R^2,\qquad
 p_0\simeq\sum_jw_je^{-\lambda_j}.
\]
The second term in the variance accounts for uncertainty in the rate arising
from uncertain breeding value. Replacing the mixture by a single Poisson
distribution with rate \(\mu_R\) would discard that contribution.

\subsection{A closed-form prediction for the quadratic log mean}

For the quadratic model, the first two rate moments can also be obtained
analytically. This provides a check independent of the likelihood
optimisation.

\begin{proposition}[Normal average of an exponential quadratic]
\label{estimation:closed_prediction}
Let \(X\sim N(M,v)\), with \(v>0\), and
\(\lambda(X)=\exp(\alpha+bX+HX^2/2)\). For \(k>0\), put
\(D_k=1-kHv\). If \(D_k>0\), then
\begin{equation}
 \E[\lambda(X)^k]
 =D_k^{-1/2}
   \exp\!\left\{
     k\alpha+kbM+\frac{kHM^2}{2}
     +\frac{k^2v(b+HM)^2}{2D_k}
   \right\}.
 \label{estimation:rate_moments}
\end{equation}
If \(D_k\leq0\), that moment is not finite. For \(v=0\), it is simply
\(\lambda(M)^k\).
\end{proposition}

\begin{proof}
Write \(X=M+\sqrt v\,T\), with \(T\sim N(0,1)\), and expand:
\[
 k\alpha+kbX+\frac{kH}{2}X^2
 =k\alpha+kbM+\frac{kHM^2}{2}
   +k\sqrt v\,(b+HM)T+\frac{kHv}{2}T^2.
\]
Multiplying by the standard normal density combines the last quadratic term
with \(-T^2/2\). Let \(d=k\sqrt v\,(b+HM)\). The remaining integral is
\[
 \frac1{\sqrt{2\pi}}\int
   \exp\!\left(-\frac{D_k}{2}t^2+dt\right)\,dt.
\]
When \(D_k>0\), complete the square:
\[
 -\frac{D_k}{2}t^2+dt
 =-\frac{D_k}{2}\left(t-\frac d{D_k}\right)^2
   +\frac{d^2}{2D_k}.
\]
The shifted normal integral contributes \(D_k^{-1/2}\), and the constant
contributes \(\exp\{d^2/(2D_k)\}\). Restoring the terms taken outside the
integral proves the formula. If \(D_k<0\), the exponent grows quadratically
in both tails. If \(D_k=0\), a nonzero linear term grows in one tail, while a
zero linear term leaves a constant integrand on the real line. Each case
has an infinite integral.
\end{proof}

Taking \(k=1\) gives the fitted mean given measurement information; taking
\(k=2\) supplies the second rate moment in
Equation~\eqref{estimation:predictive_formula}. The primary constraint
\(H\leq0\) makes both conditions automatic. If a broader model allowed
positive \(H\), a finite probability integral would not by itself establish
finite predictive moments. A Poisson probability is bounded by one, whereas
\(\lambda(X)\) and \(\lambda(X)^2\) are not.

Substituting \(M=0,v=G\) in the first-moment formula also gives population mean
output under the generating distribution:
\[
 \E[\lambda(X)]
 =\frac{e^\alpha}{\sqrt{1-HG}}
   \exp\!\left\{\frac{Gb^2}{2(1-HG)}\right\},
 \qquad 1-HG>0.
\]
This shows explicitly why the setting \(e^\alpha=4\) describes output at the
mean breeding value rather than an average over all genotypes.

\subsection{Sampling bars for observed bins}
\label{estimation:bin_variance}

The diagnostic plot sorts individuals by their observed mean measurements and
forms ten groups. Conditional on all these measurements, group membership is
fixed. For a group \(B\) containing \(n_B\) independent genotypes, define
\(\bar R_B=n_B^{-1}\sum_{i\in B}R_i\). Conditional independence gives
\begin{align}
 \E[\bar R_B\mid Y]
 &=\frac1{n_B}\sum_{i\in B}\mu_R(Y_i),\nonumber\\
 \Var(\bar R_B\mid Y)
 &=\frac1{n_B^2}\sum_{i\in B}s_R^2(Y_i).
 \label{estimation:bin_prediction}
\end{align}
The second identity follows by expanding the variance of a sum; the
cross-genotype covariance terms are zero under the model. The displayed
half-width is
\[
 1.96\,\frac{\sqrt{\sum_{i\in B}\widehat s_R^2(Y_i)}}{n_B}.
\]
The multiplier uses a normal approximation for a bin average. The calculation
conditions on fitted parameters and is displayed around the observed mean.
It is not a confidence interval for the conditional fitness curve, a
simultaneous band, or a calculation that includes fitting uncertainty. The
predicted zero fraction in a group is likewise the average of the individual
mixture probabilities, \(n_B^{-1}\sum_{i\in B}\widehat p_0(Y_i)\).

\paragraph{Biological interpretation.}
An observed mean height does not identify an exact breeding value. Comparing
its count directly with \(W_A(M_i)\) would disregard the range of breeding
values compatible with those measurements. Integrating over that range is
required both when fitting counts and when checking predictions against
observable quantities. Shared environments or related genotypes would also
introduce covariance terms omitted by the bin formula.

\section{A specific effect of substituting prediction variance}
\label{estimation:shrinkage}

The reliability calculation distinguishes biological variance \(G\) from
prediction variance. The following result isolates the consequence of
substituting the latter while keeping slope and curvature fixed on the true
breeding-value scale.

\begin{proposition}[Fixed-coefficient shrinkage substitution]
\label{estimation:shrinkage_formula}
Let \(L=\Vlin>0\), \(Q=\Vquad>0\), and replace \(\mat G\) by
\(s\mat G\), where \(0<s<1\), without changing \(\vect b\) or \(\mat H\).
Then
\begin{equation}
 L_{\rm pred}=sL,\qquad Q_{\rm pred}=s^2Q,\qquad
 \mathcal A_{g,\mathrm{pred}}=\frac{L}{L+sQ}>\frac{L}{L+Q}=\Ag.
 \label{estimation:shrinkage_share}
\end{equation}
\end{proposition}

\begin{proof}
The component formulas give
\[
 \vect b^{\mathsf T}(s\mat G)\vect b=sL,\qquad
 \tfrac12\tr\{\mat H(s\mat G)\mat H(s\mat G)\}=s^2Q.
\]
Cancel the positive \(s\) in their ratio. The difference between the new
share and the target is
\[
 \frac{L}{L+sQ}-\frac{L}{L+Q}
 =\frac{LQ(1-s)}{(L+sQ)(L+Q)}>0.
\]
\end{proof}

In the balanced scalar measurement model, \(s=R_{\rm pred}\). In the
recurring guide example, \(L=4,Q=0.125,s=0.6\), so the substituted share is
\[
 \frac{4}{4+0.6(0.125)}=0.981595\ldots,
\]
compared with \(4/(4+0.125)=0.969697\ldots\). Both absolute components have
decreased. If \(L=0<Q\), the share stays zero; if \(Q=0<L\), it stays one.
If both are zero, it is undefined. At \(s=0\), both substituted components
vanish and the ratio is undefined; the limit as positive \(s\) tends to zero
must not be mistaken for its value at zero.

The deliberately naive fit in the worked example makes a different change.
It sets the latent conditional variance to zero, refits \(\alpha,b,H\) using
the predictions as exact covariates, and substitutes their empirical
variance for \(G\). The fixed-coefficient proof therefore does not establish
the direction of its overall error. For the recorded dataset, its share
is approximately \(0.673984\), closer to \(2/3\) than the latent-value point
estimate, while both component estimates are smaller than their generating
values. Agreement of one ratio in one sample does not validate that
observation model. Heterogeneous reliabilities across traits also need not
produce a scalar covariance multiplier.

\section{What the parametric bootstrap repeats}
\label{estimation:bootstrap}

Write \(D\) for the complete observable dataset and \(T(D)\) for the complete
estimation procedure: calibrate the measurement model, fit the count model,
and calculate the components and share. The target sampling distribution is
the distribution of \(T(D)\) under repeated data generated by the specified
population model. Because its parameters are unknown, the parametric
bootstrap substitutes their fitted values \citep{Efron1979}.

For each bootstrap replicate, the supplied algorithm generates
\begin{align}
 X_i^*&\sim N(0,\widehat G),\nonumber\\
 Y_{ij}^*&=\widehat\mu+X_i^*+e_{ij}^*,
 &e_{ij}^*&\sim N(0,\widehat V_E),\nonumber\\
 R_i^*\mid X_i^*&\sim
 \operatorname{Poisson}\!\left[
 \exp\{\widehat\alpha+\widehat bX_i^*
                 +\tfrac12\widehat H(X_i^*)^2\}\right].
 \label{estimation:bootstrap_model}
\end{align}
It then calculates \(T(D^*)\), including new estimates of \(\mu,G,V_E\).
Using fixed calibrated values during these refits would omit their
sampling variation and their association with the fitted fitness
coefficients. The simulator centres new breeding values at zero relative to
\(\widehat\mu\); recalibrating the new data includes the sampling variation
in the estimated centre.

For \(N_{\rm boot}\) completed replicates, let \(T_1^*,\ldots,T_{N_{\rm boot}}^*\)
be the estimates of one quantity. Their empirical distribution is
\[
 \widehat F_{\rm boot}(t)
 =\frac1{N_{\rm boot}}\sum_{j=1}^{N_{\rm boot}}
   {\boldsymbol1}_{\{T_j^*\leq t\}}.
\]
The percentile interval uses its 0.025 and 0.975 quantiles. This specifies
how the interval is constructed. It does not prove exact 95\% coverage.
The procedure approximates a sampling distribution under the fitted model;
bias, boundaries, nonregular estimation or misspecification can affect
coverage.

There are two separate sources of variation. With the original data held
fixed, increasing \(N_{\rm boot}\) makes the simulated empirical distribution
less variable. It does not make the fitted biological or count model more
appropriate. To assess repeated-sample coverage, one would simulate many
independent original datasets, construct an interval for each, and calculate
the fraction containing the relevant generating quantity.

The default reproduction uses 500 replicates with seed 20261006, after the
observations generated with seed 20261005. All 500 recorded refits complete.
That statement concerns computational success. The code records failures.
If \(N_{\rm requested}\) fits were requested, it returns intervals only when
at least \(\max\{2,\lfloor0.9N_{\rm requested}\rfloor\}\) have completed.
The floor denotes rounding down to an integer. That numerical rule is not a
statistical justification
for ignoring failed samples: their causes and relation to the parameter
estimates must still be examined. The successful default run does not
establish interval coverage.

\paragraph{Biological interpretation.}
These intervals account for repeated sampling under the assumed additive
trait, independent measurement environments and conditional Poisson
reproduction. They do not express uncertainty that these assumptions are
biologically correct. Adding bootstrap replicates cannot account for
unmodelled relatedness, nonadditive variation, shared environments or a wrong
fitness function.

\section{A quartic alternative and the scope of the AIC comparison}
\label{estimation:aic}

The sensitivity fit extends the log mean to
\begin{equation}
 \eta_4(x)=\alpha+bx+\tfrac12Hx^2+c_3x^3+c_4x^4.
 \label{estimation:quartic_fit}
\end{equation}
It keeps the measurement calibration fixed between the two fitted families.
The quadratic family has three fitness parameters; the quartic family has
five. Applying the AIC expression \(2k-2\ell\) to their conditional
likelihoods gives
\begin{align}
 \Delta{\rm AIC}
 &={\rm AIC}_4-{\rm AIC}_2\nonumber\\
 &=(2\cdot5-2\ell_4)-(2\cdot3-2\ell_2)\nonumber\\
 &=4-2(\ell_4-\ell_2).
 \label{estimation:aic_difference}
\end{align}
A negative difference favours the more flexible fit under this criterion
\citep{Akaike1974}. The reported default difference, approximately \(3.410883\),
comes from the two numerical fitted log likelihoods.

This is an approximate conditional-likelihood comparison, not a formal test of
the whole measurement-and-reproduction model. Both fits use the same
estimated calibration. The primary family restricts \(H\leq0\); the quartic
family allows \(H\) up to three but constrains its highest coefficient to
\(-1\leq c_4\leq-10^{-8}\), with \(-2\leq c_3\leq2\).
Negative \(c_4\) makes the log mean decline in both tails and keeps
population rate moments finite. The exactly quadratic case \(c_4=c_3=0\)
lies on the limiting boundary of the quartic family, not in the interior of
these numerical bounds. Ordinary unconstrained likelihood-ratio
calibrations should not be inferred from the comparison.

\subsection{Why the cubic term is needed when the centre moves}
\label{estimation:quartic_centring}

Suppose a generating curve is written in a coordinate \(x\) as
\[
 \ell(x)=\alpha+bx+\tfrac12Hx^2-\kappa x^4.
\]
Move the reference point by \(\delta\), so \(x=z+\delta\). Expanding
\((z+\delta)^2=z^2+2\delta z+\delta^2\) and
\[
 (z+\delta)^4=z^4+4\delta z^3+6\delta^2z^2+4\delta^3z+\delta^4
\]
gives
\begin{align}
 \ell(z+\delta)
 &=\alpha_\delta+b_\delta z+\tfrac12H_\delta z^2
       -4\kappa\delta z^3-\kappa z^4,\nonumber\\
 \alpha_\delta&=\alpha+b\delta+\tfrac12H\delta^2-\kappa\delta^4,\nonumber\\
 b_\delta&=b+H\delta-4\kappa\delta^3,\qquad
 H_\delta=H-12\kappa\delta^2.
 \label{estimation:quartic_shift}
\end{align}
A quartic curve with no cubic term at one centre generally has a cubic term
at another. Including both terms allows the fitted family to represent that
change of centre. In the design used here, \(c_3\) multiplies \(z^3\) and
\(c_4\) multiplies \(z^4\) directly; the third and fourth derivatives at zero
are \(6c_3\) and \(24c_4\). The coefficient \(H\) still multiplies \(z^2/2\)
and is the second derivative at zero.

\section{Local quadratic variation and the complete quartic curve}
\label{estimation:quartic}

Let \(X\sim N(0,G)\), and consider the exact generating log mean
\begin{equation}
 \ell(X)=\alpha+bX+\tfrac12HX^2-\kappa X^4,
 \qquad \kappa\geq0.
 \label{estimation:quartic_truth}
\end{equation}
The fourth-order term has first derivative \(-4\kappa x^3\) and second
derivative \(-12\kappa x^2\), both zero at \(x=0\). Therefore its local
quadratic approximation at the population mean still has slope \(b\),
curvature \(H\), and the share derived from \(b^2G\) and \(H^2G^2/2\).
Its complete population variance requires additional moments.

\begin{proposition}[Variance and affine fit of the quartic curve]
\label{estimation:quartic_variance}
Under Equation~\eqref{estimation:quartic_truth},
\begin{align}
 \Var\{\tfrac12HX^2-\kappa X^4\}
 &=\tfrac12H^2G^2-12H\kappa G^3+96\kappa^2G^4,\nonumber\\
 R_{\rm full}^2
 &=\frac{b^2G}
 {b^2G+\tfrac12H^2G^2-12H\kappa G^3+96\kappa^2G^4},
 \label{estimation:quartic_full_r2}
\end{align}
provided the full-curve variance is positive. The second expression is the
population affine-regression \(R^2\) of the complete log mean on \(X\).
\end{proposition}

\begin{proof}
The Gaussian moment results in Chapter~\ref{gaussian:chapter} give
\[
 \E[X^2]=G,\qquad \E[X^4]=3G^2,\qquad
 \E[X^6]=15G^3,\qquad \E[X^8]=105G^4.
\]
Consequently,
\begin{align*}
 \Var(X^2)
 &=\E[X^4]-\E[X^2]^2=3G^2-G^2=2G^2,\\
 \Cov(X^2,X^4)
 &=\E[X^6]-\E[X^2]\E[X^4]
   =15G^3-3G^3=12G^3,\\
 \Var(X^4)
 &=\E[X^8]-\E[X^4]^2
   =105G^4-9G^4=96G^4.
\end{align*}
Apply the variance-of-a-sum identity, retaining the cross term:
\begin{align*}
 \Var\{\tfrac12HX^2-\kappa X^4\}
 &=\frac{H^2}{4}\Var(X^2)
   +\kappa^2\Var(X^4)
   +2\left(\frac H2\right)(-\kappa)\Cov(X^2,X^4)\\
 &=\tfrac12H^2G^2+96\kappa^2G^4-12H\kappa G^3.
\end{align*}
All centred odd Gaussian moments vanish. Therefore
\[
 \Cov(X,X^2)=0,\qquad \Cov(X,X^4)=0,
\]
and the linear term is uncorrelated with the even part.
It follows that \(\Cov\{X,\ell(X)\}=bG\).
For \(G>0\), the population least-squares slope is
\(\Cov\{X,\ell(X)\}/G=b\). Its intercept is the population mean
\[
 \E[\ell(X)]=\alpha+\tfrac12HG-3\kappa G^2.
\]
The fitted component therefore has variance \(b^2G\), and the total
variance is this component plus the even-term variance just derived.
Dividing proves Equation~\eqref{estimation:quartic_full_r2}.
If \(G=0\), the target is constant and its regression \(R^2\) is undefined.
\end{proof}

At \(b=0.7,H=-0.7,G=1,\kappa=0.08\), the local quadratic components remain
\(0.49\) and \(0.245\), so their share is \(2/3\). For the complete curve,
the two extra contributions to the even-part variance are
\[
 -12H\kappa G^3=0.672,\qquad
 96\kappa^2G^4=0.6144.
\]
Thus
\[
 R_{\rm full}^2
 =\frac{0.49}{0.49+0.245+0.672+0.6144}
 =0.242406253\ldots.
\]
The change is not a contradiction: the numerator describes the same linear
term, while the full-curve denominator includes additional population
variation absent from the local quadratic approximation.

\paragraph{Biological interpretation.}
Slope and curvature at the population mean do not determine how much higher
orders contribute over all breeding values present. A fitted quadratic can
estimate a local description accurately while describing the complete curve
poorly. Its bootstrap interval cannot correct that difference in target.
Even the local share extracted from a fitted quartic is not the \(R^2\) of
the complete fitted quartic unless the higher terms contribute no additional
relevant variation.

\section{Relating the derivations to the implementation}

The functions in
\codepath{examples/estimating_variance_share/model.py}
correspond directly to the steps above:
\codepath{calibrate} estimates the measurement model;
\codepath{fit_surface} evaluates and optimises the mixture likelihood;
\codepath{checked_surface} assesses integration stability;
\codepath{conditional_prediction} produces observable-data means, variances and
zero probabilities; \codepath{bootstrap} repeats the two-stage analysis; and
\codepath{truth} calculates the generating local share and complete quartic
\(R^2\). The batch calculation and Shiny app call these same functions.

The accompanying \codepath{test_model.py} checks normal moments, separation of
observations from simulation truth, the closed mean formula, an independent
adaptive likelihood integral, the known-value limit, repeated-seed recovery,
bootstrap recalibration, quartic variation, and explicit numerical failures.
Those checks establish reproducibility and agreement for the tested cases.
The proofs in this chapter identify the assumptions under which the formulas
hold; they do not turn those finite checks into evidence that an observation
model is suitable for every dataset.
